\chapter{الكيمياء فوق الجزيئية}\label{ch:b3:supramolecular}

في سنة 1962 عزل تشارلز بيدرسن، وهو يبحث عن شيء آخر، بولي إيثر
حلقيًا أذاب برمنغنات البوتاسيوم في البنزين وجعله
بنفسجيًا. فقد التفّت الحلقة حول أيون البوتاسيوم وأخفت
شحنته خلف سطح هيدروكربوني، ساحبةً البرمنغنات معها أيونًا
مضادًا. والكيمياء التي تتجاوز الجزيء، أي كيمياء التجمعات التي تمسكها
قوى أضعف من الروابط التساهمية، بدأت بذلك المحلول البنفسجي.
ويقيس هذا الفصل تلك القوى، ويفسّر لماذا تختار بعض المضيفات ضيوفها
اختيارًا جيدًا إلى هذا الحد، ويبيّن كيف تُقاس ثوابت الارتباط، وينتهي بجزيئات
مخترقة ومتشابكة في صورة آلات.

\begin{recall}
عالج مجلد السنة الأولى القوى بين الجزيئية، والرابطة الهيدروجينية، والإذابة،
والجزيئات الكارهة للماء، والمعقدات، وثوابت التكوين والمعقدات المخلبية؛ وعالج مجلد السنة الثانية
الجهد الكيميائي، وطاقات جيبس والتوفيق بالمربعات الصغرى.
ووصف \cref{ch:b3:advanced-nmr} التبادل السريع والبطيء والاندماج،
و\cref{ch:b3:rate-theories} \omterm{thm:b3:rate-theories:eyring}{معادلة أيرنغ}، 
و\cref{ch:b3:partition-functions} الإنتروبيا الانتقالية لغاز.
\end{recall}

\section{التآثرات غير التساهمية}

\begin{definition}[الكيمياء فوق الجزيئية]\label{def:b3:supramolecular:supramolecular}
\emph{الكيمياء فوق الجزيئية}\index{كيمياء فوق جزيئية} هي كيمياء
تجمعات الجزيئات التي تمسكها تآثرات غير تساهمية. وفي
\emph{معقد المضيف--الضيف}\index{معقد مضيف--ضيف}، يكون \emph{المضيف}\index{مضيف}
الجزيء الأكبر، الذي له تجويف أو شق ومواقع ارتباط متجهة
إليه، و\emph{الضيف}\index{ضيف} الجزيء الأصغر أو الأيون الذي يربطه.
\end{definition}

التآثرات هي تآثرات مجلد السنة الأولى، أيون--ثنائي قطب، وثنائي قطب--ثنائي قطب،
والروابط الهيدروجينية وقوى لندن، يضاف إليها بضعة تآثرات لها أسماؤها الخاصة.

\begin{definition}[التآثرات غير التساهمية]\label{def:b3:supramolecular:interactions}
\emph{تراص $\pi$}\index{تراص باي@تراص $\pi$} هو التجاذب بين حلقات عطرية مستوية
متقابلة وجهًا لوجه، مزاحة أو مائلة عادة. و\emph{تآثر الكاتيون--$\pi$}\index{تآثر كاتيون--باي@تآثر كاتيون--$\pi$}
هو التجاذب بين كاتيون 
والوجه الغني بالإلكترونات لحلقة عطرية. و\emph{الرابطة الهالوجينية}\index{رابطة هالوجينية}
هي التجاذب بين الطرف الفقير بالإلكترونات لذرة هالوجين مرتبطة
تساهميًا، المقابل لرابطتها، وقاعدة لويس. و\emph{المفعول الكاره
للماء}\index{مفعول كاره للماء} هو ميل السطوح غير القطبية إلى التلاقي
في الماء، الذي يدفعه أساسًا تحرير جزيئات الماء المنتظمة
حولها.
\end{definition}

تآثرات الأيون--ثنائي القطب والكاتيون--$\pi$ هي الأقوى، وتتوقف
كثيرًا على المذيب الذي ينافس على الأيون؛ تليها الروابط الهيدروجينية،
ثم تراص $\pi$ والروابط الهالوجينية، ثم قوى لندن، الضعيفة فرادى لكنها
تُجمع على سطوح تماس كبيرة. وفي الماء، كثيرًا ما يغلب \omterm{def:b3:supramolecular:interactions}{المفعول الكاره للماء}:
\omterm{def:b3:supramolecular:supramolecular}{فالمضيف} المصمم ليربط ضيفًا بروابط هيدروجينية في الكلوروفورم قد لا
يربطه إطلاقًا في الماء، الذي تنافس روابطه الهيدروجينية الخاصة.

\section{التعرف الجزيئي}

\begin{definition}[التعرف الجزيئي]\label{def:b3:supramolecular:recognition}
\emph{التعرف الجزيئي}\index{تعرف جزيئي} هو الارتباط
الانتقائي \omterm{def:b3:supramolecular:supramolecular}{لضيف} \omterm{def:b3:supramolecular:supramolecular}{بمضيف}. ويعتمد على \emph{التتامّ}\index{تتام}،
أي توافق الحجم والشكل ومواقع الارتباط بين \omterm{def:b3:supramolecular:supramolecular}{المضيف} \omterm{def:b3:supramolecular:supramolecular}{والضيف}، 
ويعززه \emph{التنظيم المسبق}\index{تنظيم مسبق}، أي \omterm{def:b3:supramolecular:supramolecular}{مضيف} تكون
مواقع ارتباطه ممسوكة سلفًا في هندسة المعقد قبل الارتباط.
\end{definition}

ثابت الارتباط $K_a = [\mathrm{HG}]/([\mathrm H][\mathrm G])$ ثابت تكوين بمعنى
مجلد السنة الأولى، و$\Delta G^\circ = -RT\ln K_a$. \omterm{def:b3:supramolecular:supramolecular}{والمضيف} المرن يجب أن يجمّد
تشكيله عند الارتباط، وهذا يكلّف إنتروبيا؛ أما \omterm{def:b3:supramolecular:supramolecular}{المضيف} الصلب المنظم مسبقًا فقد
دفع هذا الثمن في أثناء تخليقه.

\begin{definition}[المفعول المخلبي ومفعول الحلقة الكبيرة]\label{def:b3:supramolecular:macrocyclic}
\emph{المفعول المخلبي}\index{مفعول مخلبي} هو الاستقرار الأكبر
لمعقد مع ربيطة متعددة المخالب منه مع عدد مكافئ من
الربائط الأحادية المخلب ذات الذرات المانحة نفسها. و\emph{مفعول الحلقة
الكبيرة}\index{مفعول الحلقة الكبيرة} هو الاستقرار الإضافي حين تكون
الربيطة المتعددة المخالب حلقة لا سلسلة مفتوحة.
\end{definition}

\begin{proposition}[المفعول المخلبي إنتروبي أساسًا]\label{prop:b3:supramolecular:chelate-entropy}
في $\mathrm{[M(NH_3)_6]^{2+} + 3\,en \rightleftharpoons [M(en)_3]^{2+} + 6\,NH_3}$، حيث en هو الإيثان-1،2-ثنائي أمين، توجد
روابط فلز--نيتروجين الست نفسها في الطرفين، ويدفع التفاعلَ
ازديادُ عدد الجزيئات الحرة.
\end{proposition}

\begin{proof}[تعليل]
الروابط المتكوّنة والمنكسرة من النوع نفسه، فتكون $\Delta_rH^\circ$ صغيرة. ويزداد عدد
الجزيئات المستقلة من أربعة إلى سبعة: أي ثلاث درجات حرية انتقالية
إضافية لجزيء كامل، كل منها يساوي عشرات الجول لكل كلفن
ولكل مول في المحلول (\cref{ch:b3:partition-functions})، ومنه $\Delta_rS^\circ > 0$ و$K > 1$.
ومن الناحية الحركية، حالما يرتبط أحد طرفي الربيطة المخلبية، يُمسك الطرف الآخر قرب
الفلز بتركيز فعال مرتفع، فيرتبط على الفور إذا انفك.
\end{proof}

\begin{proposition}[التعويض بين الأنتالبي والإنتروبي]\label{prop:b3:supramolecular:compensation}
في سلسلة من معقدات المضيف--الضيف المتقاربة، تقترن قيمة $\Delta H^\circ$ أكثر سلبية للارتباط
عادة بقيمة $\Delta S^\circ$ أكثر سلبية، بحيث تتغير $\Delta G^\circ$ أقل من كل منهما (وهذه
ملاحظة تجريبية).
\end{proposition}

المعقد الأكثر إحكامًا يرتبط بقوة أكبر لكنه يقيّد حركات \omterm{def:b3:supramolecular:supramolecular}{المضيف} 
\omterm{def:b3:supramolecular:supramolecular}{والضيف} أكثر؛ والرابطة الهيدروجينية الأقوى مع \omterm{def:b3:supramolecular:supramolecular}{الضيف} تعني غالبًا رابطة أضعف مع
الماء الذي أزاحه. وطاقة جيبس، التي تحدد الانتقائية، هي الفرق الصغير
بين حدين كبيرين.

\begin{definition}[المضيفات]\label{def:b3:supramolecular:hosts}
\emph{الإيثر التاجي}\index{إيثر تاجي} بولي إيثر حلقي كبير، مثل
18-تاج-6، $\ce{(CH2CH2O)6}$. و\emph{الكريبتاند}\index{كريبتاند} بولي إيثر
ثنائي الحلقة له رأسا جسر نيتروجينيان، يحيط بكاتيون في ثلاثة
أبعاد في \emph{كريبتات}\index{كريبتات}. 
و\emph{السيكلودكسترين}\index{سيكلودكسترين} قليل وحدات حلقي من وحدات الغلوكوز، على
شكل مخروط مقطوع ذي تجويف كاره للماء ومجموعات هيدروكسي على
حافتيه. و\emph{الكاليكسارين}\index{كاليكسارين} حلقة كبيرة على شكل كأس من
وحدات فينول تصلها جسور ميثيلين.
\end{definition}

\begin{omfigure}
\begin{tikzpicture}[font=\small]
  % 18-crown-6 with K+
  \foreach \k in {0,...,17} {
    \pgfmathsetmacro\ang{90 + 20*\k}
    \pgfmathsetmacro\nxt{110 + 20*\k}
    \draw[black!70] (\ang:1.55) -- (\nxt:1.55);
  }
  \foreach \k in {0,3,...,15} {
    \pgfmathsetmacro\ang{90 + 20*\k}
    \node[fill=white, inner sep=0.5pt, text=omThm] (o\k) at (\ang:1.55) {O};
    \draw[omThm!60, densely dashed] (0,0) -- (\ang:1.32);
  }
  \node[circle, fill=omDef!25, draw=omDef, inner sep=1pt] at (0,0) {$\ce{K+}$};
  \node[font=\scriptsize] at (0,-2.1) {\foreignlanguage{arabic}{18-تاج-6 مع $\ce{K+}$}};
  % [2.2.2]cryptand with K+
  \begin{scope}[xshift=5.2cm]
    \node (nt) at (0,1.7) {N};
    \node (nb) at (0,-1.7) {N};
    \foreach \x/\n in {-1.25/a, 0/b, 1.25/c} {
      \node[text=omThm, fill=white, inner sep=0.5pt] (oa\n) at (\x,0.55) {O};
      \node[text=omThm, fill=white, inner sep=0.5pt] (ob\n) at (\x,-0.55) {O};
      \draw[black!70] (nt) to[out=-90+40*\x, in=90] (oa\n) -- (ob\n) to[out=-90, in=90-40*\x] (nb);
    }
    \node[circle, fill=omDef!25, draw=omDef, inner sep=1pt] at (0.62,0) {$\ce{K+}$};
    \node[font=\scriptsize] at (0,-2.4) {\foreignlanguage{arabic}{كريبتاند \babelsublr{[2.2.2]}: كريبتات $\ce{K+}$}};
  \end{scope}
\end{tikzpicture}

{\small يسارًا: 18-تاج-6، حلقة من اثنتي عشرة ذرة كربون (الرؤوس) وست ذرات
أكسجين، في مركزها أيون بوتاسيوم ترتبط به ذرات الأكسجين الست. يمينًا:
\omterm{def:b3:supramolecular:hosts}{الكريبتاند} [2.2.2]، ثلاث سلاسل من ذرتي أكسجين إيثريتين بين رأسي جسر
نيتروجينيين، رسم تخطيطي: يُمسك الكاتيون داخل قفص ثلاثي الأبعاد،
بست ذرات أكسجين وذرتي النيتروجين.}
\end{omfigure}

تنتقي الإيثرات التاجية الكاتيونات بحسب حجمها: فالمركب 18-تاج-6 يربط $\ce{K+}$ أقوى من
$\ce{Na+}$ الأصغر، الذي لا يستطيع بلوغ ذرات الأكسجين الست في آن واحد، ومن $\ce{Cs+}$ الأكبر، الذي
يجلس فوق الحلقة. والكريبتاندات، الأكثر تنظيمًا مسبقًا والتي تلف الأيون
كليًا، ترتبط بقوة أكبر أيضًا وتنتقي على نحو أدق. ويسحب الكاتيون، المختبئ داخل
\omterm{def:b3:supramolecular:supramolecular}{المضيف}، أنيونه إلى المذيبات غير القطبية، حيث يصير الأنيون،
غير المذاب، شديد التفاعل: وهذا أساس الحفز بنقل الطور.

\section{قياس الارتباط}

\begin{proposition}[التبادل السريع]\label{prop:b3:supramolecular:fast-exchange}
حين يتبادل \omterm{def:b3:supramolecular:supramolecular}{مضيف} بين حالتيه الحرة والمرتبطة أسرع بكثير من
الفرق بين ترددي رنينهما، تظهر إشارته في الرنين المغناطيسي النووي عند المتوسط
$\delta_{\mathrm{obs}} = x_{\mathrm{free}}\delta_{\mathrm{free}} + x_{\mathrm{bound}}\delta_{\mathrm{bound}}$، حيث $x$ نسبتا \omterm{def:b3:supramolecular:supramolecular}{المضيف} في كل حالة.
\end{proposition}

\begin{proof}
يزور كل جزيء \omterm{def:b3:supramolecular:supramolecular}{مضيف} الحالتين مرات كثيرة في أثناء القياس؛ 
وتردد مبادرته، بمتوسطه الزمني، هو المتوسط المرجَّح للترددين، بأوزان
هي نسب الزمن المقضي في كل حالة، وهي عند التوازن
نسب الجزيئات في كل حالة (\cref{ch:b3:advanced-nmr}).
\end{proof}

\begin{theorem}[متساوي الحرارة الدقيق 1:1]\label{thm:b3:supramolecular:isotherm}
من أجل $\mathrm{H + G \rightleftharpoons HG}$ بتركيزين كليين $H_0$ و$G_0$ وثابت ارتباط $K$،
\[
  [\mathrm{HG}] = \frac{1}{2}\Bigl[\Bigl(H_0 + G_0 + \frac1K\Bigr) - \sqrt{\Bigl(H_0 + G_0 + \frac1K\Bigr)^2 - 4H_0G_0}\,\Bigr].
\]
\end{theorem}

\begin{proof}
مع $c = [\mathrm{HG}]$، $K = c/((H_0 - c)(G_0 - c))$، أي $c^2 - (H_0 + G_0 + 1/K)c + H_0G_0 = 0$. ومن الجذرين،
لا يقع تحت $H_0$ و$G_0$ كليهما إلا الأصغر (فجداؤهما $H_0G_0$ ومجموعهما
يتجاوز $H_0 + G_0$): وهو ذو إشارة الناقص.
\end{proof}

\begin{omfigure}
\begin{tikzpicture}
\begin{axis}[name=a, width=0.48\linewidth, height=5.5cm, xmin=0, xmax=6, ymin=0, ymax=1.05,
  axis lines=left, xlabel={\foreignlanguage{arabic}{مكافئات الضيف، $G_0/H_0$}}, ylabel={\foreignlanguage{arabic}{نسبة المضيف المرتبطة}},
  tick label style={font=\scriptsize}, label style={font=\small},
  ]
\addplot[omMeth, thick] table[x=equiv, y=K4] {figdata/out/bachelor-3-binding-titration.dat};
\addplot[omDef, thick] table[x=equiv, y=K3] {figdata/out/bachelor-3-binding-titration.dat};
\addplot[omThm, thick] table[x=equiv, y=K2] {figdata/out/bachelor-3-binding-titration.dat};
\node[font=\scriptsize, omMeth, anchor=north] at (axis cs:4,0.94) {$K = 10^4$};
\node[font=\scriptsize, omDef, anchor=north] at (axis cs:4,0.68) {$K = 10^3$};
\node[font=\scriptsize, omThm, anchor=north] at (axis cs:4,0.27) {$K = 10^2$};
\end{axis}
\begin{axis}[at={(a.east)}, anchor=west, xshift=1.4cm, width=0.48\linewidth, height=5.5cm,
  xmin=0, xmax=1, ymin=0, ymax=1.05, axis lines=left,
  xlabel={\foreignlanguage{arabic}{الكسر المولي للمضيف}}, ylabel={$[\text{المعقد}]/c_{\text{الكلي}}$},
  tick label style={font=\scriptsize}, label style={font=\small}]
\addplot[omDef, thick] table[x=x, y=HG] {figdata/out/bachelor-3-binding-titration-b.dat};
\addplot[omThm, thick] table[x=x, y=HG2] {figdata/out/bachelor-3-binding-titration-b.dat};
\draw[black!35, dashed] (axis cs:0.5,0) -- (axis cs:0.5,0.55) (axis cs:0.3333,0) -- (axis cs:0.3333,0.38);
\node[font=\scriptsize, omDef, anchor=south] at (axis cs:0.5,0.52) {HG};
\node[font=\scriptsize, omThm, anchor=south] at (axis cs:0.25,0.36) {$\mathrm{HG_2}$};
\end{axis}
\end{tikzpicture}

{\small يسارًا: نسبة \omterm{def:b3:supramolecular:supramolecular}{المضيف} المرتبطة في أثناء معايرة عند $H_0 = \qty{1.0}{mmol/L}$ من أجل ثلاثة
ثوابت ارتباط (بوحدة $\unit{L.mol^{-1}}$): كلما كبر $K$ ازدادت حدة الانكسار عند مكافئ واحد؛
وحين $KH_0 \gg 1$ لا يقول المنحنى إلا إن $K$ كبير. يمينًا: مخططا جوب، أي
تركيز المعقد عند تركيز كلي ثابت \omterm{def:b3:supramolecular:supramolecular}{للمضيف} \omterm{def:b3:supramolecular:supramolecular}{والضيف} معًا،
وقمته عند نسبة \omterm{def:b3:supramolecular:supramolecular}{مضيف} 1/2 لمعقد 1:1 و1/3 لمعقد 1:2 (ثوابت
نموذجية).}
\end{omfigure}

\begin{method}[ثابت الارتباط من معايرة بالرنين المغناطيسي النووي]\label{met:b3:supramolecular:nmr-titration}
\begin{enumerate}
  \item أبقِ تركيز \omterm{def:b3:supramolecular:supramolecular}{المضيف} $H_0$ ثابتًا وأضف \omterm{def:b3:supramolecular:supramolecular}{الضيف}، من 0 إلى عدة
        مكافئات؛ وسجّل إشارة \omterm{def:b3:supramolecular:supramolecular}{للمضيف} عند كل نقطة (تبادل سريع).
  \item النموذج: $\delta = \delta_{\mathrm{free}} + (\delta_{\mathrm{bound}} - \delta_{\mathrm{free}})[\mathrm{HG}]/H_0$، مع $[\mathrm{HG}]$ من متساوي الحرارة الدقيق.
  \item وفّق $K$ و$\delta_{\mathrm{bound}}$ بالمربعات الصغرى غير الخطية، بتصغير مجموع مربعات
        البواقي؛ وخذ الأخطاء المعيارية من انحناء ذلك المجموع
        (مصفوفة التغاير).
  \item افحص البواقي بحثًا عن اتجاه (ستوكيومترية خاطئة) وتحقّق من أن
        النسبة المرتبطة تمتد تقريبًا من 20 إلى 80\,\% (فعندئذ تقيّد المعطيات $K$).
\end{enumerate}
\end{method}

\begin{definition}[مخطط جوب]\label{def:b3:supramolecular:job}
\emph{مخطط جوب}\index{مخطط جوب} (طريقة التغيرات المتصلة) منحنى
تركيز معقد، أو إشارة متناسبة معه، بدلالة
الكسر المولي لأحد المكوّنين، عند تركيز كلي ثابت
لكليهما.
\end{definition}

\begin{proposition}[قيمة مخطط جوب العظمى]\label{prop:b3:supramolecular:job}
من أجل معقد وحيد $\mathrm{H_nG_m}$، يبلغ \omterm{def:b3:supramolecular:job}{مخطط جوب} قيمته العظمى عند الكسر المولي \omterm{def:b3:supramolecular:supramolecular}{للمضيف}
$x = n/(n + m)$، مهما كان ثابت الارتباط.
\end{proposition}

\begin{proof}
مع التركيز الكلي $T$، $H_t = xT$، $G_t = (1 - x)T$، والمعقد $c$، والحر $[\mathrm H] = H_t - nc$،
$[\mathrm G] = G_t - mc$، و$c = \beta[\mathrm H]^n[\mathrm G]^m$. وبالاشتقاق بالنسبة إلى $x$ عند القيمة العظمى، حيث
$\dd c/\dd x = 0$: $0 = \beta\bigl(n[\mathrm H]^{n-1}[\mathrm G]^m T - m[\mathrm H]^n[\mathrm G]^{m-1}T\bigr)$، ومنه $n[\mathrm G] = m[\mathrm H]$. وبالتعويض،
$n(G_t - mc) = m(H_t - nc)$، ومنه $nG_t = mH_t$، أي $n(1 - x) = mx$: $x = n/(n + m)$.
\end{proof}

\begin{method}[إجراء مخطط جوب]\label{met:b3:supramolecular:job}
\begin{enumerate}
  \item حضّر محلولين أمّين متساويي التركيز المولي \omterm{def:b3:supramolecular:supramolecular}{للمضيف} \omterm{def:b3:supramolecular:supramolecular}{والضيف}.
  \item امزجهما بنسب $x$ من 0 إلى 1 عند حجم كلي ثابت.
  \item قِس خاصية متناسبة مع المعقد (امتصاصية مصحَّحة
        لإسهام الأنواع الحرة، أو $x\,\Delta\delta$ في الرنين المغناطيسي النووي).
  \item اقرأ الستوكيومترية من موضع القيمة العظمى.
\end{enumerate}
\end{method}

\begin{definition}[المعايرة الحرارية متساوية الحرارة]\label{def:b3:supramolecular:itc}
\emph{المعايرة الحرارية متساوية الحرارة}\index{معايرة حرارية متساوية الحرارة}
(ITC) تقيس الحرارة المنطلقة أو الممتصة عند كل حقنة من محلول
\omterm{def:b3:supramolecular:supramolecular}{الضيف} في محلول \omterm{def:b3:supramolecular:supramolecular}{المضيف} المحفوظ عند درجة حرارة ثابتة.
\end{definition}

\begin{proposition}[الحرارة لكل حقنة]\label{prop:b3:supramolecular:itc}
بإهمال التمديد، تكون حرارة الحقنة رقم $i$ في خلية حجمها $V_0$
$q_i = \Delta H^\circ V_0\bigl([\mathrm{HG}]_i - [\mathrm{HG}]_{i-1}\bigr)$.
\end{proposition}

\begin{proof}
الحرارة هي أنتالبي التفاعل مضروبة في كمية المعقد المتكوّنة في أثناء
الحقنة، وتلك الكمية هي تغير $[\mathrm{HG}]$ مضروبًا في حجم الخلية.
\end{proof}

\begin{method}[قراءة متساوي حرارة ITC]\label{met:b3:supramolecular:itc}
\begin{enumerate}
  \item ارسم الحرارة لكل مول من \omterm{def:b3:supramolecular:supramolecular}{الضيف} المحقون بدلالة النسبة المولية.
  \item ارتفاع الهضبة قبل الإشباع يعطي $\Delta H^\circ$؛ والنسبة المولية عند
        نقطة الانعطاف تعطي الستوكيومترية؛ والانحدار يعطي $K$.
  \item وفّق الثلاثة كلها بمتساوي الحرارة الدقيق؛ ثم $\Delta G^\circ = -RT\ln K$ 
        و$T\Delta S^\circ = \Delta H^\circ - \Delta G^\circ$.
\end{enumerate}
\end{method}

فتجربة ITC واحدة تعطي إذن البصمة الترموديناميكية الكاملة
للارتباط، الأنتالبي والإنتروبي، التي لا تعطيها معايرة بالرنين المغناطيسي النووي أو بمطيافية المرئي--فوق البنفسجي إلا
عبر سلسلة من درجات الحرارة.

\begin{inthelab}[معايرة بالرنين المغناطيسي النووي]
يوضع محلول \omterm{def:b3:supramolecular:supramolecular}{المضيف} في مذيب مدوتر، تركيزه نحو مليمول في اللتر،
في أنبوب رنين مغناطيسي نووي ويُسجَّل طيفه. ثم يضاف محلول مركّز
\omterm{def:b3:supramolecular:supramolecular}{للضيف} محضَّر في محلول \omterm{def:b3:supramolecular:supramolecular}{المضيف}، بحيث يبقى تركيز \omterm{def:b3:supramolecular:supramolecular}{المضيف}
ثابتًا، على دفعات صغيرة بمحقنة دقيقة؛ وبعد كل إضافة يُرج
الأنبوب، ويُترك ليتوازن عند درجة حرارة المسبار، ويُسجَّل
الطيف. ويُقرأ انزياح إشارة \omterm{def:b3:supramolecular:supramolecular}{للمضيف} تتحرك عند كل نقطة ثم
يُوفَّق.
\end{inthelab}

\section{المضيفات}

تحتوي السيكلودكسترينات، التي تصنعها الإنزيمات من النشاء، على ست أو سبع أو ثماني وحدات
غلوكوز ($\alpha$، $\beta$، $\gamma$). ويمسك تجويفها الكاره للماء الأجزاء غير القطبية من الضيوف في
الماء: فهي تذيب الأدوية، وتحمل العطور وتنزع الكوليسترول من
أغشية الخلايا. أما الكاليكسارينات والكوكوربيتوريلات ذات شكل اليقطين، وهي
حلقات كبيرة صلبة من وحدات الغلايكوليوريل ذات بوابات مبطنة بالكربونيل، فتربط الكاتيونات 
والكاتيونات العضوية بقوة كبيرة جدًا في الماء.

\begin{omfigure}
\begin{tikzpicture}[font=\small]
  \draw[black!70, fill=omDef!8] (-1.9,1.2) -- (-1.2,-1.2) arc[start angle=180, end angle=360, x radius=1.2, y radius=0.3]
    -- (1.9,1.2);
  \draw[black!70, fill=omDef!18] (0,1.2) ellipse[x radius=1.9, y radius=0.45];
  \draw[black!40, dashed] (-1.2,-1.2) arc[start angle=180, end angle=0, x radius=1.2, y radius=0.3];
  \node[draw=omThm, thick, regular polygon, regular polygon sides=6, minimum size=8mm, inner sep=0pt] at (0,0.2) {};
  \draw[omThm, thick] (0,0.62) -- (0,2.0);
  \node[font=\scriptsize, omThm, anchor=west] at (0.1,2.0) {\foreignlanguage{arabic}{الضيف}};
  \node[font=\scriptsize, anchor=west] at (2.0,1.2) {\foreignlanguage{arabic}{الحافة الواسعة: \babelsublr{OH} ثانوية}};
  \node[font=\scriptsize, anchor=west] at (1.3,-1.25) {\foreignlanguage{arabic}{الحافة الضيقة: \babelsublr{OH} أولية}};
  \node[font=\scriptsize, anchor=east] at (-1.75,0) {\foreignlanguage{arabic}{تجويف كاره للماء}};
\end{tikzpicture}

{\small \omterm{def:b3:supramolecular:hosts}{سيكلودكسترين}، مرسوم في صورة المخروط المقطوع الذي يكوّنه، مع \omterm{def:b3:supramolecular:supramolecular}{ضيف}
عطري في تجويفه. مجموعات الهيدروكسي على الحافتين تجعل الخارج
ذوابًا في الماء؛ أما الداخل، المبطن بمجموعات C--H وذرات الأكسجين الغليكوزيدية، فهو
غير قطبي.}
\end{omfigure}

\section{التجمع الذاتي والآلات الجزيئية}

\begin{definition}[التجمع الذاتي]\label{def:b3:supramolecular:self-assembly}
\emph{التجمع الذاتي}\index{تجمع ذاتي} هو الانتظام التلقائي والعكوس
لمكوّنات في بنية محددة، بتحكم من
المعلومات التي تحملها المكوّنات نفسها: أشكالها ومواقع
ارتباطها.
\end{definition}

الأيونات الفلزية ذات زوايا التناسق الثابتة والربائط الجسرية الصلبة تتجمع
في مربعات وأقفاص وكبسولات في خطوة واحدة، لأن كل رابطة يمكن أن تنفتح
وتنغلق حتى تُبلغ البنية الأكثر استقرارًا، دون إجهاد أو مواقع متدلية.
وتزاوج القواعد في الحمض النووي، وطيّ البروتينات، هما الفكرة نفسها
على مقياس أكبر.

\begin{definition}[الجزيئات المتشابكة ميكانيكيًا]\label{def:b3:supramolecular:interlocked}
\emph{الجزيء المتشابك ميكانيكيًا}\index{جزيء متشابك ميكانيكيًا}
له أجزاء غير مرتبطة بعضها ببعض لكن لا يمكن فصلها
دون كسر رابطة تساهمية. و\emph{الروتاكسان}\index{روتاكسان} حلقة
منظومة على محور ذي سدادتين ضخمتين في طرفيه؛ 
و\emph{الكاتينان}\index{كاتينان} مكوّن من حلقات متشابكة.
\end{definition}

\begin{definition}[التخليق بالقالب]\label{def:b3:supramolecular:template}
\emph{التخليق بالقالب}\index{تخليق بالقالب} يستعمل تآثرًا
مؤقتًا، مع أيون فلزي غالبًا، لإمساك مكوّنات متفاعلة في
الهندسة اللازمة لتفاعل كان سيكون بعيد الاحتمال لولا ذلك، مثل
إغلاق حلقة حول حلقة أخرى.
\end{definition}

\begin{omfigure}
\begin{tikzpicture}[font=\small, >=Stealth]
  % step 1: two crescents around Cu+
  \draw[omDef, line width=2pt] (0.3,0) arc[start angle=0, end angle=300, radius=0.75];
  \draw[omThm, line width=2pt] (-0.3,0) arc[start angle=180, end angle=-120, radius=0.75];
  \fill[omProp] (0,0) circle (0.13);
  \node[font=\scriptsize] at (0,-1.35) {\foreignlanguage{arabic}{$\ce{Cu+}$ يمسك ربيطتين}};
  \node[font=\scriptsize] at (0,-1.7) {\foreignlanguage{arabic}{مفتوحتين متصالبتين}};
  \draw[->, thick] (1.4,0) -- (2.3,0) node[midway, above, font=\scriptsize] {\foreignlanguage{arabic}{إغلاق}};
  % step 2: closed interlocked rings with Cu
  \begin{scope}[xshift=3.6cm]
    \draw[omDef, line width=2pt] (-0.35,0) circle (0.75);
    \draw[omThm, line width=2pt] (0.35,0) circle (0.75);
    \draw[white, line width=5pt] ([shift={(52:0.75)}]-0.35,0) arc[start angle=52, end angle=72, radius=0.75];
    \draw[omDef, line width=2pt] ([shift={(48:0.75)}]-0.35,0) arc[start angle=48, end angle=76, radius=0.75];
    \fill[omProp] (0,0) circle (0.13);
    \node[font=\scriptsize] at (0,-1.35) {\foreignlanguage{arabic}{كاتينات}};
  \end{scope}
  \draw[->, thick] (5.0,0) -- (5.9,0) node[midway, above, font=\scriptsize] {$-\ce{Cu+}$};
  \begin{scope}[xshift=7.2cm]
    \draw[omDef, line width=2pt] (-0.35,0) circle (0.75);
    \draw[omThm, line width=2pt] (0.35,0) circle (0.75);
    \draw[white, line width=5pt] ([shift={(52:0.75)}]-0.35,0) arc[start angle=52, end angle=72, radius=0.75];
    \draw[omDef, line width=2pt] ([shift={(48:0.75)}]-0.35,0) arc[start angle=48, end angle=76, radius=0.75];
    \node[font=\scriptsize] at (0,-1.35) {\foreignlanguage{arabic}{كاتينان \babelsublr{[2]}}};
  \end{scope}
  % rotaxane
  \begin{scope}[xshift=10.2cm]
    \draw[black!70, line width=1.6pt] (-1.1,0) -- (1.1,0);
    \fill[black!60] (-1.25,0) circle (0.2) (1.25,0) circle (0.2);
    \draw[omThm, line width=2pt] (0,0) ellipse[x radius=0.18, y radius=0.55];
    \node[font=\scriptsize] at (0,-1.35) {\foreignlanguage{arabic}{روتاكسان \babelsublr{[2]}}};
  \end{scope}
\end{tikzpicture}

{\small \omterm{def:b3:supramolecular:template}{التخليق بالقالب} \omterm{def:b3:supramolecular:interlocked}{لكاتينان} (رسم تخطيطي): يمسك أيون النحاس(I) ربيطتين
مبنيتين على الفينانترولين متصالبتين بزاوية قائمة في تناسقه رباعي
الأوجه؛ وإغلاق كل منهما في حلقة يعطي حلقتين متشابكتين حول
الفلز، ونزع النحاس يترك \omterm{def:b3:supramolecular:interlocked}{الكاتينان} الحر. يمينًا: \omterm{def:b3:supramolecular:interlocked}{روتاكسان}، أي
حلقة محتجزة على محور بسدادتين.}
\end{omfigure}

\begin{definition}[الآلات الجزيئية]\label{def:b3:supramolecular:machine}
\emph{الآلة الجزيئية}\index{آلة جزيئية} تجمّع تتحرك
مكوّناته بعضها نسبةً إلى بعض حركة متحكمًا فيها استجابةً لمنبّه
(الضوء، أو تغير في الأكسدة والاختزال، أو تغير في pH)، ويمكن جعله يؤدي شغلًا.
\end{definition}

في المكوك الجزيئي، تجلس حلقة على محور \omterm{def:b3:supramolecular:interlocked}{روتاكسان} عند إحدى محطتين؛ 
والمنبّه الذي يضعف ارتباطها بالأولى يرسلها إلى الثانية، 
والمنبّه المعاكس يعيدها. والمحركات الدوّارة المدفوعة بالضوء المبنية على
ألكينات شديدة الازدحام تدور في اتجاه واحد فقط، بتناوب تماكبات
ضوئية وخطوات حرارية لا يمكن أن تجري عكسيًا. وتُقاس سرعة
التنقل بين المحطتين بالرنين المغناطيسي النووي متغير درجة الحرارة، من
اندماج إشارتي الصورتين (\cref{ch:b3:advanced-nmr}).

\begin{safety}{\ghs[0.62cm]{GHS07}}
% ledger: ghs4:18C6
18-تاج-6: ضار إذا ابتُلع، مهيج للجلد والعينين والمسالك التنفسية.
والإيثرات التاجية تحمل الأملاح الفلزية إلى المذيبات العضوية وعبر الجلد؛ فهي
توزن وتُتداول بالقفازات.
\end{safety}

\begin{history}[الكيمياء فوق الجزيئية وآلاتها]
أسست الإيثراتُ التاجية لتشارلز بيدرسن (1967)، وكريبتانداتُ جان ماري لين ومضيفاتُ دونالد
كرام المنظمةُ مسبقًا هذا الميدانَ؛ وتقاسم الثلاثة جائزة نوبل في
الكيمياء لسنة 1987. وصنع جان بيير سوفاج الكاتينانات بمردود مرتفع بقالب
النحاس سنة 1983، وبنى فريزر ستودارت مكوكات \omterm{def:b3:supramolecular:interlocked}{الروتاكسان}، وبنى بن
فيرينغا محركًا دوّارًا مدفوعًا بالضوء سنة 1999؛ وتقاسموا جائزة نوبل في
الكيمياء لسنة 2016 عن تصميم \omterm{def:b3:supramolecular:machine}{الآلات الجزيئية} وتخليقها.
\end{history}

\section{تمارين}

\begin{exercise}[$\star$]\label{exo:b3:supramolecular:1}
سمِّ التآثر الرئيس في: $\ce{K+}$ في 18-تاج-6؛ حلقتي بنزين متقابلتين وجهًا لوجه؛
$\ce{K+}$ فوق حلقة بنزين؛ يودوبنتافلوروبنزين مع البيريدين؛ زوج
غوانين--سيتوزين.
\end{exercise}

\begin{exercise}[$\star$]\label{exo:b3:supramolecular:2}
احسب $\Delta G^\circ$ عند \qty{298}{K} \omterm{def:b3:supramolecular:supramolecular}{لمعقد مضيف--ضيف} فيه $K_a = \qty{1.0e4}{L.mol^{-1}}$.
\end{exercise}

\begin{exercise}[$\star$]\label{exo:b3:supramolecular:3}
في الميثانول، $\log K = 6.1$ من أجل $\ce{K+}$ و4.4 من أجل $\ce{Na+}$ مع 18-تاج-6 (معطيات التمرين).
احسب الانتقائية والفرق في $\Delta G^\circ$.
\end{exercise}

\begin{exercise}[$\star$]\label{exo:b3:supramolecular:4}
يبلغ \omterm{def:b3:supramolecular:job}{مخطط جوب} قمته عند كسر مولي \omterm{def:b3:supramolecular:supramolecular}{للمضيف} قدره 0.33. ما الستوكيومترية؟
\end{exercise}

\begin{exercise}[$\star\star$]\label{exo:b3:supramolecular:5}
يتلقى \omterm{def:b3:supramolecular:supramolecular}{مضيف} تركيزه \qty{2.0}{mmol/L} مع $K = \qty{1500}{L.mol^{-1}}$ و$\delta_{\mathrm{free}} = \qty{3.560}{ppm}$ و$\delta_{\mathrm{bound}} = \qty{3.720}{ppm}$ مكافئًا
واحدًا من \omterm{def:b3:supramolecular:supramolecular}{الضيف}. احسب النسبة المرتبطة والانزياح الملاحظ.
\end{exercise}

\begin{exercise}[$\star\star$]\label{exo:b3:supramolecular:6}
من أجل النيكل(II)، $\log\beta_6 = 8.6$ مع الأمونيا و$\log\beta_3 = 18.3$ مع الإيثان-1،2-ثنائي أمين (معطيات
التمرين). احسب الثابت و$\Delta G^\circ$ عند \qty{298}{K} للتفاعل
$\ce{[Ni(NH3)6]^2+} + 3\,\mathrm{en} \rightleftharpoons \ce{[Ni(en)3]^2+} + \ce{6NH3}$، وقل ما الذي يدفعه.
\end{exercise}

\begin{exercise}[$\star\star$]\label{exo:b3:supramolecular:7}
تعطي تجربة ITC عند \qty{298}{K} القيمتين $K = \qty{2.0e5}{L.mol^{-1}}$ و$\Delta H^\circ = \qty{-40}{kJ/mol}$. احسب $\Delta G^\circ$ و$T\Delta S^\circ$ 
و$\Delta S^\circ$.
\end{exercise}

\begin{exercise}[$\star\star$]\label{exo:b3:supramolecular:8}
فسّر دور النحاس(I)، ودور هندسته الرباعية الأوجه، في تخليق سوفاج
\omterm{def:b3:supramolecular:interlocked}{للكاتينان}. وكيف يُنزع النحاس في النهاية؟
\end{exercise}

\begin{exercise}[$\star\star$]\label{exo:b3:supramolecular:9}
تندمج إشارتا مكوك \omterm{def:b3:supramolecular:interlocked}{روتاكسان}، المتباعدتان بمقدار \qty{120}{Hz}، عند \qty{300}{K} (معطيات
التمرين). احسب ثابت سرعة التبادل عند الاندماج وطاقة جيبس
لتنشيط التنقل.
\end{exercise}

\begin{exercise}[$\star\star\star$]\label{exo:b3:supramolecular:10}
لدواء ذوبانية ذاتية قدرها \qty{0.10}{mmol/L} ويكوّن معقدًا 1:1 مع
\omterm{def:b3:supramolecular:hosts}{سيكلودكسترين}، $K = \qty{500}{L.mol^{-1}}$ (معطيات التمرين). بيّن أن الذوبانية الكلية في
محلول تركيزه الكلي من \omterm{def:b3:supramolecular:hosts}{السيكلودكسترين} $C_t$ هي $S_0 + KS_0C_t/(1 + KS_0)$، واحسبها من أجل $C_t = \qty{10}{mmol/L}$.
\end{exercise}

\begin{exercise}[$\star\star\star$]\label{exo:b3:supramolecular:11}
بيّن انطلاقًا من متساوي الحرارة الدقيق أنه حين $KH_0 \gg 1$ تساوي النسبة المرتبطة
عدد المكافئات حتى الواحد، ثم تبقى عند الواحد. ولماذا لا تعطي مثل هذه المعايرة
إلا حدًا أدنى لقيمة $K$؟
\end{exercise}

\begin{exercise}[$\star\star\star$]\label{exo:b3:supramolecular:12}
من أجل \omterm{def:b3:supramolecular:supramolecular}{مضيف} \cref{exo:b3:supramolecular:5}، افترض أن الارتباط
\omterm{def:b3:rate-theories:diffusion}{محكوم بالانتشار}، $k_{\mathrm{on}} = \qty{1e9}{L.mol^{-1}.s^{-1}}$. احسب $k_{\mathrm{off}}$ وبيّن أن التبادل سريع على
المقياس الزمني للرنين المغناطيسي النووي عند \qty{400}{MHz}.
\end{exercise}

\section{مسألة: البنزين البنفسجي}

\begin{problem}[{مسألة نهاية الأسبوع --- البنزين البنفسجي: البوتاسيوم في إيثر تاجي،
واستخلاص البرمنغنات إلى مذيب عضوي، وثابت ارتباط
من معايرة بالرنين المغناطيسي النووي، ولماذا تكفي كمية حفزية من الإيثر التاجي}]
\label{pb:b3:supramolecular:1}
% ledger: const:R, ghs:KMnO4, ghs:toluene, rion:K+VI, rion:Na+VI
معطيات المسألة. في الميثانول عند \qty{298}{K}، $\log K = 6.1$ من أجل $\ce{K+}$ و4.4 من أجل $\ce{Na+}$ مع
18-تاج-6. الاستخلاص: يُرج طور مائي فيه \qty{0.10}{mol/L} من $\ce{KCl}$ و\qty{1.0}{mmol/L} من $\ce{KMnO4}$
مع حجم مساوٍ من التولوين يحوي \omterm{def:b3:supramolecular:hosts}{الإيثر التاجي} L؛ والنوع الوحيد المستخلص
هو الزوج الأيوني $\ce{[KL]+}\ce{MnO4-}$، مع
$K_{\mathrm{ex}} = [\ce{[KL]+MnO4-}]_{\mathrm{org}}/([\ce{K+}]_{\mathrm{aq}}[\ce{MnO4-}]_{\mathrm{aq}}[\mathrm L]_{\mathrm{org}}) = \qty{1.0e5}{L^2.mol^{-2}}$، ويبقى \omterm{def:b3:supramolecular:hosts}{الإيثر التاجي} في التولوين. ومعايرة بالرنين المغناطيسي النووي
\omterm{def:b3:supramolecular:hosts}{لإيثر تاجي} (\qty{2.0}{mmol/L}) بملح للصوديوم في الميثانول المدوتر، إشارة $\ce{CH2}$ واحدة
(ppm) بدلالة مكافئات $\ce{Na+}$:

\begin{center}\small
\begin{tabular}{l*{11}{c}}
\hline
مكافئات & 0 & 0.25 & 0.5 & 0.75 & 1.0 & 1.5 & 2.0 & 3.0 & 4.0 & 6.0 & 8.0\\
$\delta$ & 3.560 & 3.590 & 3.612 & 3.635 & 3.649 & 3.674 & 3.688 & 3.697 & 3.704 & 3.708 & 3.714\\
\hline
\end{tabular}
\end{center}
% table: binding-titration-c

\medskip\noindent\textbf{الجزء الأول --- المعقد.}
\begin{enumerate}
  \item ارسم 18-تاج-6 ومعقده مع البوتاسيوم.
  \item احسب $\Delta G^\circ$ لارتباط $\ce{K+}$ ولارتباط $\ce{Na+}$.
  \item احسب الانتقائية $\ce{K+}/\ce{Na+}$.
  \item مع أنصاف الأقطار للتناسق السداسي $r(\ce{K+}) = \qty{138}{pm}$ و$r(\ce{Na+}) = \qty{102}{pm}$، فسّر الانتقائية.
  \item لماذا يذوب المعقد في التولوين بينما لا يذوب $\ce{KCl}$؟
  \item لماذا يرتبط \omterm{def:b3:supramolecular:hosts}{الإيثر التاجي} بقوة أقل في الماء منه في الميثانول؟
\end{enumerate}

\medskip\noindent\textbf{الجزء الثاني --- الاستخلاص.}
\begin{enumerate}[resume]
  \item اكتب توازن الاستخلاص.
  \item مع $y$ التركيز المستخلص، اكتب المعادلة التي يحققها $y$ حين
        $[\mathrm L]_0 = \qty{1.0}{mmol/L}$.
  \item حلّها، وأعطِ نسبة البرمنغنات المستخلصة.
  \item احسب النسبة مع $[\mathrm L]_0 = \qty{10}{mmol/L}$.
  \item احسبها مع $[\mathrm L]_0 = \qty{0.10}{mmol/L}$.
  \item لماذا يكون التولوين بنفسجيًا، ولماذا تكون البرمنغنات المستخلصة مؤكسدًا أقوى
        منها في الماء؟
\end{enumerate}

\medskip\noindent\textbf{الجزء الثالث --- معايرة بالرنين المغناطيسي النووي.}
\begin{enumerate}[resume]
  \item لماذا تتحرك إشارة متوسطة واحدة بدل إشارتين؟
  \item اكتب النموذج الذي يعطي $\delta$ بدلالة تركيز \omterm{def:b3:supramolecular:supramolecular}{الضيف}.
  \item وفّق $K$ و$\delta_{\mathrm{bound}}$ بالمربعات الصغرى؛ وأعطِ $K$ مع خطئه المعياري.
  \item أعطِ تشتت البواقي وعلّق.
  \item على أي مجال من المكافئات تكون النسبة المرتبطة بين 20 و80\,\%؟
  \item لماذا تخفق المعايرة نفسها في قياس $K$ للبوتاسيوم، $\log K = 6.1$؟
\end{enumerate}

\medskip\noindent\textbf{الجزء الرابع --- الحفز بنقل الطور.}
\begin{enumerate}[resume]
  \item يتأكسد ألكين في التولوين بالبرمنغنات المستخلصة. أين يجري
        التفاعل؟
  \item ماذا يحدث \omterm{def:b3:supramolecular:hosts}{للإيثر التاجي} بعد التفاعل، ولماذا تكفي كمية
        حفزية منه؟
  \item لماذا يفيد فائض $\ce{KCl}$ في الماء؟
  \item ما الحفازات الأرخص التي تؤدي العمل نفسه في الصناعة؟
  \item لماذا يُستعمل التولوين اليوم حيث استعمل بيدرسن البنزين؟
  \item اذكر النتيجة: نسبة البرمنغنات المستخلصة عند
        $[\mathrm L]_0 = \qty{1.0}{mmol/L}$.
\end{enumerate}
\end{problem}
